\documentclass[10pt]{article}

% UCL dissertation draft in the official TMLR layout. For a genuinely
% anonymous TMLR submission, remove [preprint]. Do not modify the style files.
\usepackage[preprint]{vendor/tmlr}

% Stable project notation used throughout the manuscript.
\newcommand{\Bzero}{B0}
\newcommand{\Bone}{B1}
\newcommand{\BtwoM}{B2-M}
\newcommand{\BtwoD}{B2-D}
\newcommand{\Tone}{T1}
\newcommand{\Ttwo}{T2}
\newcommand{\CARLA}{CARLA}
\newcommand{\NLL}{\mathrm{NLL}}
\newcommand{\ADE}{\mathrm{ADE}}
\newcommand{\FDE}{\mathrm{FDE}}
\newcommand{\risk}{\alpha}
\newcommand{\model}{f_{\theta}}


\usepackage{amsmath,amssymb}
\usepackage{booktabs}
\usepackage{graphicx}
\usepackage{array}
\usepackage{xcolor}
\usepackage{hyperref}
\usepackage{url}

\graphicspath{%
  {../../paper/generated/distinction_v1/11_w1_manuscript/}%
  {../../paper/generated/distinction_v1/08_corrected_closed_loop/r3_final/synthesis/}%
  {../../paper/generated/paper_assets_v1/figures/}%
}
\hypersetup{colorlinks=true,linkcolor=blue,citecolor=blue,urlcolor=blue}
\newcolumntype{L}[1]{>{\raggedright\arraybackslash}p{#1}}

\title{Task-Adapted Motion Prediction under\\
Predictor--Risk Coupling: A Controlled CARLA\\
Give-Way Study}

\author{\name UCL MSc Candidate \\
  \addr Candidate identity and programme metadata withheld in the W1 review draft \\
  University College London}

\begin{document}
\maketitle
% UCL's assessment brief remains authoritative for the permitted-use category.
% This disclosure also satisfies TMLR's first-page transparency requirement.
\footnotetext{Generative-AI disclosure: OpenAI Codex was used as an assistive
tool for research planning, code generation and debugging, audit automation,
and manuscript drafting and editing. The candidate is responsible for
verifying the sources, code, results, interpretations and submitted text; the
tool is neither an author nor an evidence source.}

\begin{abstract}
Motion prediction and risk allocation are commonly evaluated in isolation,
although their benefit in autonomous driving is realised through a coupled
prediction--planning--supervision stack. This dissertation studies a
controlled give-way intersection in CARLA and asks whether task adaptation,
explicit interaction-sequence modelling and adaptive risk allocation improve
offline prediction and closed-loop safety--efficiency outcomes. A
rollout-disjoint dataset of 200 factorial rollouts is used to compare a
pretrained MultiPath control, a simple task-adapted prediction head,
capacity-matched multilayer perceptron controls and two Transformer adapters
under a validation-frozen, one-shot test protocol. Simple task adaptation
% EVIDENCE: H1_B1_MINUS_B0_TEST_ADE, H1_B1_MINUS_B0_TEST_FDE,
% H1_B1_MINUS_B0_TEST_NLL
% SOURCE: supervisor_feedback_v1/03_finetune_audit/FINETUNE_AUDIT_MANIFEST.json
reduces held-out rollout-macro uncalibrated NLL from 2.171 to
1.857~nats/step, ADE from 1.283 to 0.100~m and FDE from 2.644 to 0.121~m
relative to the pretrained control, while the tested
Transformer variants use their sequence input but do not outperform the
simple adapter. A prospectively frozen corrected closed-loop matrix contains
80 rollouts spanning two predictor stacks, four risk policies, two target
styles and five new paired initialisation groups. The prediction manipulation
favours the adapted stack in all 40/40 metric--policy--style checks, yet only
% EVIDENCE: H3_SUPPORTED_CELLS, H3_PRESPECIFIED_CELLS
% SOURCE: distinction_v1/10_four_hypothesis_evidence/M1_EVIDENCE_MANIFEST.json
2/8 closed-loop cells jointly improve completion and separation. Adaptive
% EVIDENCE: H4_DOMINANCE_CELLS, H4_PRESPECIFIED_CELLS
% SOURCE: distinction_v1/10_four_hypothesis_evidence/M1_EVIDENCE_MANIFEST.json
risk dominates a fixed comparator in only 3/12 prespecified comparisons.
Thus task adaptation, rather than added interaction-model complexity, provides
the directionally consistent in-distribution offline gain in this setting,
while its downstream value is conditional on the risk--control context.
\end{abstract}

\section{Introduction}
\label{sec:introduction}

\subsection{Motivation and context}

Autonomous vehicles approaching an unsignalised give-way junction must reason
about another road user's future motion while trading progress against risk.
In a modular stack, the realised manoeuvre is not determined by the predictor
alone: its distribution is consumed by a risk-allocation rule and trajectory
solver, may be overridden by a safety supervisor, and changes the next
observation. Consequently, a model can improve displacement error or
likelihood without producing a uniform closed-loop benefit.

This dissertation began with the planning-level hypothesis that adaptive risk
allocation would outperform fixed risk in the give-way scenario. Pilot
analysis showed that a shared supervisor and predictor could compress the
observable difference between planning policies. Rather than treating this as
a failed comparison, the final study follows the mechanism upstream: it first
tests whether scenario-specific prediction can be improved, then whether
explicit temporal interaction modelling adds value, and finally whether the
frozen predictor improvement changes the closed-loop safety--efficiency
frontier.

\subsection{Research gap}

The dissertation addresses three linked gaps. First, motion-prediction studies
often stop at open-loop ADE, FDE or likelihood, leaving the downstream value of
an improvement uncertain. Second, architectural complexity and attention are
frequently compared without a simple task-adaptation baseline and an
approximately capacity-matched non-attention control. Third, a claim that an
adaptive risk rule is superior can be an artefact of comparing it with a
single fixed setting rather than multiple prespecified fixed operating points
and of omitting moderators such as predictor and target behaviour.

\TODO{Support each gap with a critical synthesis of primary literature. State
what each cited paper evaluates, what it leaves unresolved, and why the three
gaps intersect in this system.}

\paragraph{Central thesis.}
In this controlled give-way setting, task adaptation rather than additional
interaction-model complexity provides the reliable offline prediction gain;
after that gain passes through calibration, risk allocation, optimisation and
supervision, its closed-loop value becomes context dependent. Consequently,
neither better offline prediction nor adaptive risk alone establishes uniform
system-level superiority.

\subsection{Research questions and hypotheses}

The study asks:

\begin{description}
  \item[RQ1] Does task adaptation improve in-distribution give-way trajectory
  prediction relative to a pretrained MultiPath control?
  \item[RQ2] Do explicit interaction sequences and Transformer encoders add
  benefit beyond simple adaptation and approximately matched MLP capacity?
  \item[RQ3] Do offline prediction gains translate consistently into
  closed-loop benefit across risk policies and target styles?
  \item[RQ4] Does adaptive risk allocation dominate the fixed-risk
  safety--efficiency frontier?
\end{description}

Each research question maps to one core hypothesis: task adaptation improves
prediction relative to the pretrained control (H1); explicit-sequence
Transformer adaptation adds benefit beyond simple task adaptation that is not
explained by matched non-attention capacity (H2);
the offline gain transfers consistently to closed-loop safety--efficiency
outcomes (H3); and adaptive risk dominates the fixed-risk frontier (H4).
Sequence ablation, calibration-tail analysis, deployment gates and
collision-filtered retraining are mechanism or robustness checks, not
additional headline hypotheses. Table~\ref{tab:hypotheses} reports the four
evidence-bounded verdicts.

\subsection{Contributions}

The dissertation makes four focused contributions:

\begin{enumerate}
  \item evidence that simple task adaptation, rather than the tested increase
  in interaction-model complexity, is the reliable source of offline gain;
  \item mechanism diagnostics showing that sequence use and aggregate
  calibration quality are insufficient to establish model superiority;
  \item a frozen closed-loop factorial evaluation demonstrating that predictor
  value is conditional on risk policy and target style; and
  \item a frontier-based re-evaluation of the original adaptive-risk claim,
  supported by rollout-disjoint data, frozen selection and auditable
  robustness checks.
\end{enumerate}

The contribution is not a claim of general Transformer superiority or
real-road safety. It is an empirical account of where task adaptation helps,
where added sequence complexity does not, and why predictor quality and risk
allocation must be assessed as a coupled system.

\subsection{Paper organisation}

Section~\ref{sec:related} critically reviews prediction, interaction modelling
and risk-aware planning. Section~\ref{sec:problem} defines the coupled problem.
Sections~\ref{sec:methods} and~\ref{sec:experiments} specify the reproducible
protocol, Section~\ref{sec:results} reports the evidence, and
Sections~\ref{sec:discussion}--\ref{sec:conclusion} interpret the findings and
their boundaries.

\section{Related work}
\label{sec:related}

\subsection{Risk-aware autonomous intersection planning}

Chance-constrained planning converts a probabilistic forecast into a bound on
constraint violation, but the interpretation of that bound depends on the
uncertainty model and on whether the mathematical assumptions remain valid in
deployment. Iterative risk allocation was introduced to redistribute a joint
chance budget toward active constraints rather than impose an equally
conservative allocation everywhere \citep{ono2008risk}. In driving, robust
scenario MPC represents uncertain multimodal obstacles with multiple scenarios
\citep{batkovic2021scenario}; Branch MPC instead plans a trajectory tree over
reactive modes and uses a coherent risk measure to trade performance against
robustness \citep{chen2022branch}. Both illustrate that ``risk'' is an
optimisation design variable, not the same quantity as measured collision or
physical clearance.

The closest control lineage to this dissertation is the SMPC formulation of
\citet{nair2022smpc}, which embeds Gaussian-mixture predictions in multimodal
chance constraints and evaluates feedback-policy variants at a CARLA
intersection. Its later extension jointly optimises feedback policies and
mode-wise risk allocation and demonstrates real-time closed-loop operation in
several scenarios \citep{nair2025predictive}. Complementary work establishes
recursive-feasibility and safety properties for nominal, robust and
contingency chance-constrained planners under explicit Gaussian-mixture
assumptions \citep{ren2025recursively}. These are stronger formal and
cross-scenario contributions than the empirical study here; however, a formal
property of the planned policy does not by itself quantify executed
bounding-box separation after solver failure handling and safety supervision.

The present work therefore asks a narrower empirical question. Adaptive risk
is compared with aggressive, medium and conservative fixed operating points,
not with a single convenient baseline. Predictor stack and target style are
crossed with policy, and executed completion, physical separation, binary
failures, solver identity and supervisor activity are retained. This design
addresses RQ4 while avoiding the stronger claim that a finite simulator matrix
validates a chance constraint or proves real-world safety.

\subsection{Multimodal prediction, scoring and adaptation}

MultiPath represents future motion by a discrete distribution over trajectory
anchors with regressed offsets and uncertainty \citep{chai2020multipath}. This
compact Gaussian-mixture output is particularly suitable for an optimisation
consumer. MultiPath++ subsequently replaced dense raster-only fusion with
sparse scene elements, learned anchors and output aggregation
\citep{varadarajan2022multipathpp}. The dissertation retains the original
MultiPath decoder and frozen anchors so that adaptation changes a deployed
component without replacing the planner interface.

Displacement and probabilistic scores answer different questions. Top-1 ADE
and FDE assess one selected trajectory and can ignore whether probability mass
is distributed appropriately across plausible modes. Negative log likelihood
uses the full predictive density and is a proper score under the relevant
distributional conditions \citep{gneiting2007scoring}, but a low value can
arise through both better means and changed covariance. Calibration is again
distinct: neural-network confidence can be systematically miscalibrated
\citep{guo2017calibration}, and post-hoc calibrated regression can improve
coverage without changing the underlying point predictor
\citep{kuleshov2018calibrated}. Subgroup calibration matters because an
aggregate guarantee need not hold on identifiable subpopulations
\citep{hebertjohnson2018multicalibration}; trajectory-specific work likewise
finds that uncertainty estimates require calibration before planning
\citep{cao2024cctr}. These distinctions motivate reporting ADE, FDE,
uncalibrated NLL, coverage and calibrated NLL, with the response-active tail
kept separate from the aggregate.

Adaptation is also a distribution question rather than merely an architecture
choice. UniTraj reports substantial cross-dataset degradation and improvement
from greater data diversity \citep{feng2024unitraj}. Here, the intended claim
is deliberately in-distribution: rollouts are disjoint by initialisation but
come from one Town05 give-way process. Retraining the final prediction head is
therefore a task-adaptation control for domain mismatch, not evidence of
cross-map generalisation.

\subsection{Interaction-aware and Transformer prediction}

Self-attention provides content-dependent interaction across a sequence
\citep{vaswani2017attention}. Driving predictors have used it at substantially
greater scale and with richer representations than this dissertation.
AgentFormer jointly attends across time and agents
\citep{yuan2021agentformer}; Scene Transformer supports marginal, joint and
ego-conditioned multi-agent prediction in a unified architecture
\citep{ngiam2022scene}; AutoBots models socially consistent joint futures as a
sequence of agent sets \citep{girgis2022autobots}. Wayformer studies early,
late and hierarchical fusion of heterogeneous scene modalities
\citep{nayakanti2023wayformer}, while MTR couples global intention queries with
local motion refinement \citep{shi2022mtr}. These results establish attention
as a strong design family, not a default optimum for every data regime.

Architecture comparisons are confounded when representation, residual output,
parameter count or optimisation budget changes simultaneously. The current
study therefore uses two MLP/Transformer pairs with the same input and
residual-output scope and approximately similar adapter capacity. It also
retains a simpler, larger task-adapted head and reports every trainable count
and boundary-epoch selection. Frozen-input zeroing and shuffling test whether
the Transformer actually uses its six-token history, but cannot prove semantic
interaction reasoning. RQ2 is consequently about the tested adapters under a
fixed small-data budget; it is not a referendum on the larger models above.

\subsection{Prediction-aware planning and closed-loop evaluation}

Prediction and planning can be coupled during training or evaluation.
Topological multimodal prediction has been demonstrated with an
optimisation-based navigator at uncontrolled CARLA intersections
\citep{roh2021topological}. DIPP back-propagates through a differentiable
planner and reports both open- and closed-loop gains
\citep{huang2023dipp}. The nuPlan benchmark similarly argues that open-loop
distance metrics are insufficient for long-horizon planning and supplies
reactive closed-loop evaluation \citep{caesar2021nuplan}. These systems cover
broader planning tasks than the controlled junction here, but they also change
more of the stack; they do not isolate how a frozen probabilistic predictor
interacts with a risk-policy frontier.

Task relevance provides the key evaluation principle. \citet{farid2023task}
propagate prediction failures to planning cost rather than treating every
forecast error as equally harmful. \citet{nakamura2025regret} use closed-loop
regret to identify prediction failures that degrade the robot, and find that a
small high-regret subset can be informative for fine-tuning. Most directly,
the recent preprint by \citet{bouzidi2025closing} reports that open-loop model
ranking need not predict closed-loop performance and that smaller predictors
can match larger ones in its tested stack. The present dissertation agrees with
that broad warning, but contributes a different controlled contrast: it first
verifies that B1 remains better than B0 in-loop, then tests translation within
each predictor--risk--target cell and records solver/supervisor mechanisms.

\subsection{Synthesis and study positioning}

Across the reviewed literature, four claims form a hierarchy rather than a
single model leaderboard: adaptation can improve a forecast; an interaction
architecture can add value beyond matched controls; an offline improvement can
benefit closed-loop execution; and adaptive risk can dominate fixed operating
points. Prior work supplies strong methods for each layer, but the directly
comparable studies reviewed above do not evaluate all four with a frozen
predictor interface, rollout-disjoint selection, several fixed-risk controls,
paired target styles and executed mechanism telemetry. This dissertation fills
that bounded empirical gap. Its contribution is a controlled claim audit in a
single give-way system, not a new universal predictor or a formal safety proof.

\section{Problem formulation}
\label{sec:problem}

\subsection{Coupled prediction--planning system}

At time $t$, the predictor $\model$ consumes a raster observation $r_t$ and,
for sequence models, an interaction history $s_{t-H:t}$ with validity mask
$m_{t-H:t}$. It outputs a $K$-component trajectory distribution
\begin{equation}
  p_{\theta}(Y_{t+1:t+T}\mid r_t,s_{t-H:t},m_{t-H:t})
  = \sum_{k=1}^{K}\pi_k\,
  \mathcal{N}(Y;\mu_k,\Sigma_k).
\end{equation}
A risk policy maps this distribution and the current state to a risk budget
$\risk_t$. A trajectory solver then constructs a candidate ego trajectory;
the supervisor may override it before execution. Figure~\ref{fig:chain}
summarises this implemented feedback path.

\begin{figure}[t]
  \centering
  \includegraphics[width=0.98\linewidth]{figure08_deployment_chain.png}
  \caption{Deployment chain and measurement boundaries. Prediction, risk
  allocation, optimisation, supervision and execution form a feedback system.
  The arrows show implemented information flow, not causal identification.}
  \label{fig:chain}
\end{figure}

\subsection{Offline objectives}

The primary selection metric is rollout-macro mixture NLL per valid prediction
step. Secondary metrics are top-1 ADE and FDE. Macro aggregation first averages
valid windows within a rollout and then averages rollouts, avoiding the
incorrect treatment of overlapping windows as independent experiments.

\subsection{Closed-loop objectives}

Closed-loop performance is multi-objective. The corrected primary efficiency
outcome is event-clock completion time. The primary continuous safety-margin
outcome is minimum vehicle-footprint separation computed from actual actor
bounding boxes. Native/footprint collision, fixed-geometry yield failure and
completion failure are binary guards. Solver failure and supervisor activity
are mechanism measurements, not interchangeable safety outcomes. A policy
dominates another only if completion and separation are no worse and no excess
binary scientific failure occurs within the declared comparison scope.

\subsection{Scope of inference}

The target population is the controlled Town05 give-way geometry under the
specified routes and target styles. It excludes other maps, junctions,
weather, traffic populations and real vehicles. Corrected R3 uses five newly
generated paired initialisation groups (101--105) as the independent clusters
for uncertainty calculations.

\section{Methodology}
\label{sec:methods}

\subsection{CARLA scenario and system architecture}

Experiments use a controlled give-way junction in CARLA's Town05
\citep{dosovitskiy2017carla}. The ego follows a fixed route and encounters an
assertive or reactive target vehicle. Fixed-risk and adaptive-risk policies
share the trajectory solver and safety supervisor, making the executed motion
a coupled system outcome. \TODO{Insert CARLA version, synchronous step,
vehicle blueprints, route coordinates, weather, solver settings, supervisor
thresholds and A3 authority-regime definition from the frozen configuration.}

\subsection{Data collection and representation}

The factorial collection contains 200 rollouts: 50 ego initialisations crossed
with two target styles and two collection policies. Each prediction sample
contains a deployment-equivalent raster and six interaction tokens covering
$-1.0$ to $0.0$~s. Each token has 12 relative position, velocity, yaw and
distance features. The target future covers ten steps at 0.2~s intervals, for
a 2.0~s horizon, in the current target-local right-handed frame. Target style,
policy, factorial cell, split and initialisation ID are prohibited inputs.

Rollouts, not windows, are split: initialisations 1--40 train, 41--45
validation and 46--50 test. Normalisation is fit on train only. Samples with no
future label are excluded; partial-horizon samples use a masked loss.

\begin{figure}[t]
  \centering
  \includegraphics[width=0.98\linewidth]{figure01_research_workflow.png}
  \caption{Research workflow and evidence gates. The study progresses from
  the original planning question through controlled data, matched model
  comparison and frozen closed-loop evaluation; mixed and negative findings
  remain part of the evidence chain.}
  \label{fig:workflow}
\end{figure}

\begin{table}[t]
  \caption{Rollout-disjoint dataset split and usable prediction windows}
  \label{tab:dataset}
  \centering
  \begin{tabular}{lrrrrr}
    \toprule
    Split & Init IDs & Rollouts & Raw & Usable & Full horizon \\
    \midrule
    Train & 1--40 & 160 & 9,121 & 4,036 & 2,596 \\
    Validation & 41--45 & 20 & 1,034 & 506 & 326 \\
    Test & 46--50 & 20 & 1,075 & 495 & 315 \\
    \bottomrule
  \end{tabular}
\end{table}

\subsection{Prediction models}

\Bzero{} is the frozen pretrained MultiPath control. \Bone{} retrains only the
final MultiPath dense prediction head. \BtwoM{} and \BtwoD{} add flattened-
sequence MLP adapters; the former updates trajectory means and the latter
updates means, covariance/orientation and mixture logits. \Tone{} and \Ttwo{}
replace the corresponding MLP adapter with a masked Transformer. The
\BtwoM{}/\Tone{} and \BtwoD{}/\Ttwo{} parameter gaps are 10.48\% and 5.89\%,
below the preregistered 20\% matched-capacity ceiling.

\begin{figure}[t]
  \centering
  \includegraphics[width=0.98\linewidth]{figure02_model_architecture.png}
  \caption{Matched architecture and capacity controls. \Bone{} is the simple
  adaptation control; the MLP/Transformer pairs isolate attention from
  approximate adapter capacity while sharing the base prediction contract.}
  \label{fig:architecture}
\end{figure}

\TODO{Add an exact hyperparameter table: layer widths, heads, embedding size,
dropout, optimiser, learning rate, batch size, loss terms, epoch/early-stopping
rule, anchor count, frozen/trainable parameters and software versions.}

\subsection{Training, selection and calibration}

Five candidate variants are trained with seeds 11, 23 and 37 under the same
split, anchors, decoder, loss, optimiser budget and validation ranking metric.
The variant is selected by median validation rollout-macro NLL and a
representative seed is chosen using validation only. Selection is frozen before
the test set is opened. Context zero/shuffle ablations and collision-filtered
sensitivity runs are post-selection diagnostics and do not alter the primary
model selection.

Calibration parameters are fit on validation and frozen for test/deployment.
The closed-loop contrast uses the identity-calibrated \Bzero{} stack and the
validation-calibrated \Bone{} stack; it must therefore be interpreted as a
frozen predictor-stack treatment rather than a weights-only intervention.

\subsection{Statistical analysis}

Effect estimates can average the balanced paired conditions, but uncertainty
and exact sign-flip tests operate on five ego-initialisation cluster means.
The smallest attainable non-zero two-sided exact p-value is therefore 0.0625.
Holm correction is applied within each predefined family. Results emphasise
effect sizes, cluster intervals, sign consistency and exact/Holm p-values;
windows and simulator steps are never treated as independent replications.

\subsection{Data-quality audit and sensitivity}

Logging recorded 253 collision callbacks in six reactive training rollouts.
Episode reconstruction reduced these to 16 contact episodes and 62 unique
frames involving the target and infrastructure, not 253 ego--target crashes.
Because old callbacks cannot be aligned exactly to the sample clock, all 162
usable windows from the six rollouts are conservatively removed in a
validation-only sensitivity analysis. The primary dataset remains frozen; the
sensitivity asks only whether the architecture ranking changes.

\section{Experimental design}
\label{sec:experiments}

\subsection{Offline model-selection experiment}

The selection experiment is $5$ variants $\times$ $3$ seeds, giving 15
validation runs. Primary comparison: rollout-macro NLL. Secondary comparisons:
top-1 ADE/FDE, reactive-subset ADE and latency. A one-shot frozen test evaluates
the validation-selected checkpoint of each variant and the pretrained \Bzero{}
control. This tests H1 and H3 without test-set tuning.

\subsection{Sequence-use and calibration diagnostics}

For \Tone{} and \Ttwo{}, valid interaction tokens are replaced by zeros or
shuffled while weights remain frozen. An increase in NLL supports input use,
but not causal interaction understanding. Aggregate and response-active
calibration are reported separately. These analyses test H2 and diagnose
tail behaviour; they cannot reselect the model.

\subsection{Deployment reliability gate}

Day9 runs eight smoke arms through raster/sequence input, GMM decoding,
calibration, risk allocation, solver and supervisor. Passing this gate supports
deployability but is not an effect experiment. Formal closed-loop conclusions
are based only on Days10--12.

\subsection{Nominal closed-loop factorial}

The nominal design is
\[
2\ \text{predictors}\times4\ \text{risk policies}\times
2\ \text{target styles}\times5\ \text{initialisations}=80
\]
rollouts. Predictors are \Bzero{} and frozen \Bone{}. Risk policies are fixed
aggressive, fixed medium, fixed conservative and adaptive. The matrix evaluates
the full fixed frontier rather than a single hand-picked control.

\subsection{Arrival-timing robustness and pooled synthesis}

The shifted-timing design is
\[
2\ \text{predictors}\times2\ \text{risk policies}\times
2\ \text{target styles}\times2\ \text{offsets}\times
5\ \text{initialisations}=80
\]
rollouts, using fixed medium/adaptive risk and target offsets $-3$ and $+3$~m.
Combining these with matching nominal ($0$~m) conditions yields 24 cells and
120 formal rollouts for the three-offset synthesis. This design tests H4--H6
and records delay, footprint separation, solver failure, supervisor activity
and observed reliability events.

\subsection{Controls and validity safeguards}

Key controls are the pretrained \Bzero{}, simple \Bone{}, matched MLP adapters,
fixed-risk frontier, target-style factorial, paired initialisations, train-only
normalisation, validation-only selection, frozen calibration/deployment hashes,
and a collision-filtered sensitivity. Table~\ref{tab:threats} later separates
each mitigation from the limitation that remains.


\section{Results}
\label{sec:results}

Results are reported in the same H1--H4 order as the experimental design. All
200 planned prediction-data rollouts completed. Raw and merged manifests agree
on 11,230 samples; the 160/20/20 rollout split has disjoint groups, no duplicate
keys and no normalisation leakage. The largest descriptive split-covariate
standardised mean difference is 0.208. The formal closed-loop experiment also
completed all 80/80 prespecified rollouts and passed its treatment, integrity,
stopping and archive checks.

\subsection{H1: task adaptation improves in-distribution prediction}

% EVIDENCE: H1_B1_MINUS_B0_TEST_NLL, H1_B1_MINUS_B0_TEST_ADE,
% H1_B1_MINUS_B0_TEST_FDE
% SOURCE: supervisor_feedback_v1/03_finetune_audit/FINETUNE_AUDIT_MANIFEST.json
Across the three validation seeds, median rollout-macro NLL ranks \Bone{}
first, followed by \BtwoD{}, \Ttwo{}, \Tone{} and \BtwoM{}. On the one-shot
frozen test, \Bone{} reduces NLL from 2.171 to 1.857~nats/step, top-1 ADE from
1.283 to 0.100~m and FDE from 2.644 to 0.121~m relative to pretrained \Bzero{}.
The corresponding rollout-macro differences are $-0.314$~nats/step,
$-1.183$~m and $-2.523$~m. All three metrics favour \Bone{} in every one of
the five held-out initialisation groups; at that independent-unit level, the
mean paired effects are $-0.316$~nats/step NLL, $-1.189$~m ADE and
$-2.544$~m FDE. The two-sided exact sign-flip value reaches its attainable
minimum of $p=0.0625$, so the evidence is directionally consistent but is not
reported as conventional $p<0.05$ significance. \Bone{} also improves ADE and
FDE relative to constant velocity, clipped constant acceleration and the
train-mean prior in all five groups. H1 is therefore supported for the frozen
Town05 give-way distribution.

\begin{table}[t]
\centering\small
\caption{Frozen B0/B1 test results at one common rollout-macro aggregation. Every metric is first averaged within rollout and then over 20 rollouts. The 315 overlapping windows are descriptive observations, not independent replications; lower is better.}
\label{tab:finetune-b0-b1-rollout-macro}
\begin{tabular}{@{}lrrrrrr@{}}
\toprule
Stack/contrast & Windows & Rollouts & Init groups & NLL & ADE (m) & FDE (m) \\
\midrule
B0 & 315 & 20 & 5 & 2.171 & 1.283 & 2.644 \\
B1 & 315 & 20 & 5 & 1.857 & 0.100 & 0.121 \\
\midrule
B1$-$B0 & 315 & 20 & 5 & -0.314 & -1.183 & -2.523 \\
\bottomrule
\end{tabular}
\end{table}


The mechanism checks narrow this conclusion. Raster shuffling increases
\Bone{} ADE by approximately 0.270--0.293~m across seeds, whereas target-history
shuffling changes aggregate ADE by about 0.00008~m on average; the aggregate
gain is therefore raster-dominant. In a small response-active diagnostic,
uncalibrated \Bone{} NLL is 2.076 versus 2.472 for \Bzero{}, but the
validation-frozen calibration changes these values to 8.573 and 2.959. Because
this subset contains only 15 overlapping windows from six rollouts and three
initialisation groups, it identifies a possible calibration boundary rather
than a confirmatory failure rate.

An earlier progress report described an improvement from $0.98\%$ to $100\%$.
That percentage measured how often the top-probability mode matched the
oracle-best mode; it was neither a trajectory-accuracy threshold nor a valid
independent-sample estimate. It is therefore withdrawn. The frozen-test
NLL/ADE/FDE comparison above, using common rollout-macro aggregation and five
paired held-out initialisations, is the evidence used for H1.

\begin{center}
  \refstepcounter{figure}
  \label{fig:offline}
  \centering
  \includegraphics[width=0.98\linewidth]{figure03_offline_model_comparison.png}
  \smallskip
  {\small Figure~\thefigure: Frozen offline model comparison. Points show validation seeds and
  horizontal segments their medians; test ADE/FDE include the pretrained
  \Bzero{} control. NLL, ADE and FDE use rollout-macro aggregation, so
  overlapping windows are not counted as independent replications.}
\end{center}

\subsection{H2: the tested Transformer adapters add no consistent value}

% EVIDENCE: H2_T1_TEST_NLL, H2_B2_M_TEST_NLL,
% H2_T2_TEST_NLL, H2_B2_D_TEST_NLL
% SOURCE: distinction_v1/10_four_hypothesis_evidence/M1_EVIDENCE_MANIFEST.json
The two matched-scope comparisons point in opposite directions. For the
mean-only scope, \Tone{} reaches a frozen-test NLL of 2.003746 compared with
2.024409 for the MLP control \BtwoM{}, an improvement of 0.020663. For the
distributional scope, \Ttwo{} reaches 1.878148 compared with 1.872789 for
\BtwoD{}, a degradation of 0.005359. Neither Transformer outperforms the
simple \Bone{} result of 1.857094. H2 is therefore not supported under the
tested architecture and optimisation budget.

Zeroing or shuffling the valid interaction sequence worsens both Transformer
variants, showing that they use the temporal input; sequence neglect does not
explain the negative result. The correct interpretation is nevertheless
bounded. \Bone{} has approximately 1.034 million trainable parameters, while
the adapters contain approximately 77--176 thousand, and 10/15 training runs
select a checkpoint at the 20-epoch ceiling. The experiment therefore supports
simple adaptation as the strongest tested configuration, not a general claim
that Transformers cannot improve interaction prediction.

\subsection{H3: offline gain does not transfer consistently to closed loop}

The manipulation check succeeds before closed-loop outcomes are interpreted.
Across 40/40 prespecified metric--policy--style checks, in-loop predictions
favour \Bone{} over \Bzero{} in all five paired initialisation groups. The
\Bone{}--\Bzero{} top-1 ADE effects range from $-1.352$ to $-1.125$~m and FDE
effects from $-3.152$ to $-2.315$~m. Thus the weak translation cannot be
explained by the two deployed predictor stacks becoming indistinguishable.

Only two of the eight risk-policy--target-style cells jointly show faster mean
completion and no-worse mean physical separation for \Bone{} relative to
\Bzero{}: adaptive/reactive and fixed-aggressive/assertive. The remaining six
fail at least one registered direction. For example, adaptive/assertive is
0.240~s slower with 0.0067~m lower separation, whereas fixed-aggressive/
assertive is 0.790~s faster with 0.0053~m higher separation. All Holm-adjusted
exact tests are non-confirmatory. H3 is therefore not supported as a consistent
offline-to-closed-loop translation claim, although the two favourable cells
show that the adapted stack can have context-specific closed-loop value.

\begin{center}
  \refstepcounter{table}
  \label{tab:h3-translation}
  \small
  Table~\thetable: Predictor-stack translation in the formal closed-loop experiment.
  A supported direction requires lower mean completion time and no-worse mean
  separation for \Bone{}--\Bzero{}. Each cell contains five paired
  initialisation groups.\par\smallskip
  \begin{tabular}{@{}llrrl@{}}
    \toprule
    Policy & Style & Completion (s) & Separation (m) & Direction \\
    \midrule
    Adaptive & Assertive &  0.240 & $-0.0067$ & Not supported \\
    Adaptive & Reactive & $-0.180$ &  0.0006 & Supported \\
    Fixed aggressive & Assertive & $-0.790$ &  0.0053 & Supported \\
    Fixed aggressive & Reactive & $-0.010$ & $-0.0106$ & Not supported \\
    Fixed conservative & Assertive & $-0.520$ & $-0.0262$ & Not supported \\
    Fixed conservative & Reactive &  0.490 & $-0.0056$ & Not supported \\
    Fixed medium & Assertive &  0.230 &  0.0072 & Not supported \\
    Fixed medium & Reactive & $-0.220$ & $-0.0084$ & Not supported \\
    \bottomrule
\end{tabular}
\end{center}

Figure~\ref{fig:h3-translation} visualises the registered joint-direction
criterion: a favourable cell must improve completion without sacrificing the
measured physical margin, rather than improving either outcome alone.

\begin{center}
  \refstepcounter{figure}
  \label{fig:h3-translation}
  \centering
  \includegraphics[width=0.99\linewidth]{figure_r3_h3_translation.png}
  \smallskip
  {\small Figure~\thefigure: Predictor-stack translation across risk policies
  and target styles. All values are derived from the formal 80-rollout
  closed-loop analysis.}
\end{center}

All adaptive cells use the adaptive solver on every audited step, every one of
the 20 adaptive rollouts varies its risk allocation within rollout, and the
fixed settings remain stable across initialisations. All 80 rollouts complete,
with no observed native collision, zero-margin footprint collision,
fixed-geometry yield failure or completion failure. These zero counts establish
nominal reliability in the tested sample; they do not prove zero population
risk, mathematical MPC feasibility or equivalence between treatments.

\subsection{H4: adaptive risk is context dependent rather than dominant}

% EVIDENCE: H4_DOMINANCE_CELLS, H4_PRESPECIFIED_CELLS
% SOURCE: distinction_v1/10_four_hypothesis_evidence/M1_EVIDENCE_MANIFEST.json
Adaptive risk satisfies the prespecified empirical dominance rule in only
3/12 predictor--style--fixed-comparator contrasts: \Bzero{}/assertive versus
fixed medium, and \Bone{}/reactive versus fixed aggressive and fixed medium.
It does not dominate in the remaining 9/12 contrasts. H4 is therefore not
supported as a universal frontier-dominance claim. The positive comparisons
retain a narrower conclusion: adaptive risk can be a useful context-dependent
operating point, but it is not a general replacement for fixed risk.

\begin{center}
  \refstepcounter{figure}
  \label{fig:h4-dominance}
  \centering
  \includegraphics[width=0.99\linewidth]{figure_r3_h4_dominance.png}
  \smallskip
  {\small Figure~\thefigure: Adaptive-risk dominance across predictor stack,
  target style and fixed comparator. Dominance is the prespecified empirical
  Pareto rule with binary guards, not an equivalence test.}
\end{center}

Solver and supervisor telemetry limits the mechanism claim. Historical
logger-unaccepted rows comprise bypass/no-solve and attempted-but-nonaccepted
contexts and cannot be relabelled wholesale as infeasible MPC optimisations.
Finite CasADi solve-time and controller-acceptance rates therefore use factual
solve attempts as their denominator, while bypass requests, raw return status
and rollout outcome remain separate. Earlier timing diagnostics show that
physical margin, controller acceptance/fallback and supervisor activity can
change together; they do not identify the supervisor as the sole cause of the
H4 pattern.

% Provisional tracked insertion. The final-only evidence builder atomically
% replaces this file after SF1--SF4 closure receipts and their source hashes pass.
\subsection{Supervisor-feedback closure analyses}
\label{sec:supervisor-feedback-closure-results}

The final behaviour, finite solver-time, controller-acceptance, downstream-
outcome and corrected-supervisor application-authority results are intentionally unavailable
in this non-submission draft.  The release audit fails closed until the immutable
R3 raw-log analyses and the preregistered SF4 matrix have passed their scientific
and provenance gates.


\begin{samepage}
\begin{center}
  \refstepcounter{table}
  \label{tab:hypotheses}
  \small
  Table~\thetable: Evidence-bounded verdicts for the four headline
  hypotheses. Counts are prespecified cells or comparisons, not simulation
  steps.\par\smallskip
  \begin{tabular}{@{}L{0.06\linewidth}L{0.46\linewidth}L{0.20\linewidth}L{0.18\linewidth}@{}}
    \toprule
    ID & Hypothesis & Verdict & Boundary \\
    \midrule
    H1 & Task adaptation improves in-distribution prediction & Supported with boundary & Frozen Town05 distribution \\
    H2 & Transformer adapters add consistent value & Not supported & Tested configurations/budget \\
    H3 & Offline gain transfers consistently to closed loop & 2/8 cells & Universal claim not supported \\
    H4 & Adaptive risk dominates fixed controls & 3/12 comparisons & Universal dominance not supported \\
    \bottomrule
\end{tabular}
\end{center}
\end{samepage}

\section{Discussion}
\label{sec:discussion}

\subsection{Why simple adaptation wins in this setting}

The evidence favours task adaptation, not a general rejection of
Transformers. The controlled distribution is narrow, the training split
contains 160 rollouts and 4,036 usable windows, the pretrained MultiPath
representation is strong, and the explicit history covers only one second. A
retrained output head can therefore correct substantial domain mismatch
without learning a new interaction representation. The sequence ablations
show that \Tone{} and \Ttwo{} do use their added input; their lower rank is not
explained by complete input neglect. Conversely, the B1 shuffles indicate that
its aggregate gain is raster-dominant rather than evidence of effective target
history modelling.

This interpretation is compatible with two strands of prior evidence. Large
interaction Transformers obtain their gains with richer multi-agent/map inputs
and much larger benchmark distributions
\citep{yuan2021agentformer,ngiam2022scene,nayakanti2023wayformer,shi2022mtr},
whereas cross-dataset analysis shows that domain transfer itself can cause a
large accuracy loss \citep{feng2024unitraj}. The present result adds the
boundary that, for a narrow in-distribution adaptation with a fixed MultiPath
interface, correcting the final head can be more valuable than adding a small
sequence adapter.

The comparison also has an important fairness boundary. The MLP/Transformer
pairs approximately match adapter capacity within two output scopes, but B1
has substantially more trainable parameters and 10/15 runs select a checkpoint
at the epoch ceiling. The result therefore supports simple adaptation as the
best tested configuration under the frozen budget, not an architecture-level
causal claim.

\subsection{Why offline gain does not imply uniform closed-loop gain}

R3 first establishes that the deployed manipulation is real: B1 improves all
40/40 in-loop prediction checks. The failure of uniform H3 translation cannot
therefore be dismissed as a missing predictor difference. The planner consumes
a calibrated distribution rather than ADE alone; risk allocation transforms
that distribution into constraints, controller acceptance/fallback and active
constraints limit candidate motion, and the supervisor can change execution. Because B1 and B0 also use
different frozen calibration, the valid estimand is the deployed predictor
stack rather than neural weights alone.

The result agrees with the task-relevance argument that only some forecast
errors alter downstream cost \citep{farid2023task,nakamura2025regret} and with
closed-loop benchmarks that reject open-loop distance as a sufficient planning
metric \citep{caesar2021nuplan}. It also independently reproduces the direction
of the recent, broader open-loop/closed-loop mismatch reported by
\citet{bouzidi2025closing}. Unlike integrated end-to-end training such as DIPP
\citep{huang2023dipp}, however, this experiment preserves the modular predictor
interface and attributes the observed effect only to the frozen stack-level
treatment.

The corrected result is deliberately modest: two of eight cells move both
completion and separation in the preferred direction, and the pattern depends
on policy and target style. Legacy response-tail and timing analyses supply
plausible mechanisms---calibration failure and simultaneous changes in solver,
supervisor and physical margin---but are not pooled with R3 and do not identify
a single causal component.

\subsection{Adaptive risk remains useful without dominance}

H4's negative universal result does not make adaptive risk valueless. It
dominates a fixed comparator in three prespecified contexts and fails to do so
in nine. Its defensible interpretation is therefore a context-dependent
operating point. Comparing it with aggressive, medium and conservative fixed
settings prevents a favourable single comparator from being presented as
general superiority. The result directly refines the project's original
planning hypothesis while retaining a substantive role for adaptive policy.

This is consistent with risk allocation as a mechanism for reducing unnecessary
conservatism rather than a theorem of universal empirical dominance
\citep{ono2008risk,nair2025predictive}. Formal SMPC safety and feasibility
results depend on their modelling assumptions \citep{ren2025recursively}; the
finite CARLA comparison instead asks which executed operating point is preferred
under the shared corrected controller and supervisor.

\subsection{What corrected-supervisor application authority explains}

% Provisional tracked insertion. The final-only evidence builder atomically
% replaces this file after every scientific and provenance gate passes.
\subsection{Supervisor-authority interpretation pending scientific closure}

The bounded interpretation of supervisor authority is intentionally unavailable
in this non-submission draft. This placeholder contains no outcome estimate,
count, percentage, probability value, effect direction, ranking or superiority
claim. The release audit remains fail-closed until the preregistered authority
matrix, source hashes and paper-integration build have all passed.


\subsection{Threats to validity}

\begin{table}[t]
  \caption{Threats, mitigations and remaining boundaries}
  \label{tab:threats}
  \centering
  \small
  \begin{tabular}{@{}L{0.25\linewidth}L{0.31\linewidth}L{0.34\linewidth}@{}}
    \toprule
    Threat & Mitigation & Remaining boundary \\
    \midrule
    315 overlapping offline windows & Rollout-macro points and five init-level paired effects & Effective inferential unit is five test init groups; minimum exact $p=0.0625$ \\
    Five corrected-R3 init groups & Paired effects, cluster intervals and exact tests & Minimum two-sided exact $p=0.0625$ \\
    Single map and junction & Predictor, policy and target-style factorial & No cross-map or real-road claim \\
    Configuration fairness & Two matched adapter scopes and full parameter/epoch audit & No pure architecture causality \\
    Predictor calibration differs & Stack-level estimand declared before R3 & No weight-only B1 effect \\
    Coupled solver/supervisor & Preregistered corrected-supervisor application-authority on/off SF4 intervention & The DID covers adaptive versus the criticised fixed-medium pair only; estimator-to-risk path and SMPC constraints remain in both arms \\
    Solver bypass and acceptance & Raw classification; factual attempts alone define finite CasADi timing and controller acceptance; bypass and raw status are separate & Acceptance is not mathematical feasibility; CasADi time is not full-stack deployment latency \\
    Zero observed binary failures & Native callbacks, zero-margin physical overlap and 0.25-m-per-actor margin-adjusted separation are reported separately & No zero-risk or equivalence claim \\
    Very small active tail & Separate labelled legacy diagnostic & 15 windows/six rollouts/three groups; no stable tail estimate \\
    Legacy implementation differs & Corrected R3 is primary; no pooling & Legacy evidence is mechanism-only \\
    \bottomrule
  \end{tabular}
\end{table}

The fixed policies are three prespecified operating points rather than a dense
continuous Pareto frontier. The implementation diagram represents information
flow, not a causal graph. These boundaries prevent overclaiming but do not
invalidate the frozen in-distribution and corrected R3 comparisons.

\subsection{Implications and future work}

For research practice, a simple task-adaptation control and training-budget
audit should precede claims for a more elaborate interaction architecture.
Closed-loop evaluation should verify that the prediction manipulation persists
in deployment and should retain policy/context cells rather than collapse them
into one score. For system design, continuous geometry margins and mechanism
telemetry remain informative when nominal binary failure counts are zero.

Future work should replicate across maps and junctions, increase independent
initialisation clusters, enrich response-active data, compare calibration-
matched predictor stacks, and evaluate larger graph/attention models with
training budgets commensurate with their capacity. These are external-validity
extensions, not missing experiments required to support the present claims.

For production-oriented evaluation, the immediate implication is procedural:
retain a cheap task-adapted control, audit probabilistic calibration by
interaction subgroup, verify that a model difference survives in-loop, and
report a policy frontier with execution mechanisms. The present sample cannot
set deployment thresholds, but this sequence reduces the risk of promoting a
larger model or an adaptive policy on an offline score alone.

\section{Conclusion}
\label{sec:conclusion}

This dissertation asked a single system-level question: when motion prediction
and risk-aware planning are coupled, which apparent component improvements
remain visible in the executed give-way behaviour? The answer is more
conditional than an isolated prediction or planning benchmark would suggest.
Simple task adaptation produces a large and directionally consistent
in-distribution prediction gain, but neither added Transformer complexity nor
adaptive risk establishes uniform system-level superiority. The main
contribution is therefore an evidence chain that follows a model change from
the frozen offline test through calibration, risk allocation, optimisation,
supervision and executed geometry.

\paragraph{H1: supported within the frozen distribution.}
Retraining the MultiPath prediction head reduces rollout-macro NLL from 2.171
to 1.857~nats/step, ADE from 1.283 to 0.100~m and FDE from 2.644 to 0.121~m
relative to the pretrained control. The direction is preserved across all five
held-out initialisation groups and against three physical baselines. This
supports task adaptation for the controlled Town05 give-way distribution, not
generalisation to other maps, interactions or real vehicles.

\paragraph{H2: not supported for the tested Transformer configurations.}
The mean-only Transformer slightly improves upon its matched MLP adapter, while
the distributional Transformer is slightly worse than its matched MLP adapter;
neither exceeds simple task adaptation. Sequence ablations confirm that the
Transformers use their temporal input, so the result is not simple input
neglect. Because the full models are not parameter matched and many runs reach
the epoch ceiling, the conclusion concerns the tested adapters and budget,
not the general value of Transformer architectures.

\paragraph{H3: not supported as a consistent transfer claim.}
The adapted predictor remains better in every in-loop manipulation check, yet
only 2/8 risk-policy--target-style cells jointly improve completion and
physical separation. Better prediction is therefore real but does not have a
uniform driving effect once it enters the coupled stack. Predictor quality
must be evaluated in the planning contexts through which it can change action.

\paragraph{H4: not supported as a universal dominance claim.}
Adaptive risk dominates a fixed comparator in 3/12 prespecified contrasts and
fails to dominate in the other nine. It should therefore be interpreted as a
context-dependent operating point rather than a universal replacement for
fixed risk. The supervisor, solver path and controller acceptance remain part
of this system-level comparison; mechanism telemetry constrains attribution
but does not justify assigning the result to the supervisor alone.

% Provisional tracked Conclusion insertion. The final-only evidence builder atomically
% replaces this file only after every scientific and provenance gate passes.
\paragraph{Supervisor-authority conclusion pending scientific closure.}
The outcome-specific supervisor-authority conclusion is intentionally unavailable
in this non-submission draft. This placeholder contains no outcome estimate,
effect direction, interval interpretation, ranking, superiority statement or
claim that the authority study has reached scientific closure.


Taken together, the four hypotheses refine the project's original expectation
that adaptive risk would simply outperform fixed risk. The stronger conclusion
is not that prediction or adaptive risk is ineffective, but that their benefit
depends on where they alter the closed-loop safety--efficiency frontier. Within
the stated simulator population, the dissertation supplies reproducible
evidence for that coupling and for the importance of retaining negative and
mixed results. It does not claim real-world safety, zero collision risk,
mathematical MPC feasibility or architecture-wide Transformer inferiority.


\subsubsection*{Broader impact statement}
This work evaluates autonomous-driving components only in a controlled CARLA
scenario and must not be interpreted as evidence of real-road safety. A
failure mode identified here is particularly relevant to deployment: global
calibration may improve aggregate likelihood while worsening rare
response-active cases. Reporting coupled solver and supervisor behaviour,
negative findings and population limits reduces the risk that an offline
metric improvement is presented as a safety guarantee. The released evidence
contains simulator trajectories and metrics rather than personal data.

\subsubsection*{Reproducibility statement}
The repository freezes dataset splits, model/configuration hashes, deployment
contracts, completion manifests, analysis tables and figure provenance. The
main experimental unit and all post-selection analyses are declared in
Sections~\ref{sec:methods}--\ref{sec:experiments}; additional audit detail is
provided in Appendices~\ref{app:reproducibility}--\ref{app:evidence}.

\bibliography{references}
\bibliographystyle{vendor/tmlr}

\appendix
\section{Reproducibility specification}
\label{app:reproducibility}

\subsection{Environment and frozen artefacts}

\TODO{Insert a table containing CARLA, Python, TensorFlow, CUDA/cuDNN, solver
and operating-system versions; hardware; Git commit; dataset manifest hash;
model-package hashes; calibration hash; random seeds; and execution commands.}

\subsection{Model and optimisation hyperparameters}

\TODO{Transcribe the exact frozen hyperparameters and trainable-parameter
counts. Do not infer missing values from memory; cite the configuration and
manifest paths in source comments.}

\subsection{Closed-loop contract}

The corrected contract contains 16 treatment cells and five paired
initialisation groups (101--105), for 80 formal rollouts. The final appendix
must list each cell identifier, fixed-risk value, adaptive-risk rule,
termination condition, infrastructure-retry policy, event-clock completion
formula and actual-bounding-box separation formula from the frozen R3/M0
artifacts. Earlier arrival-offset experiments are documented separately as
legacy diagnostics and are not part of this contract.

\subsection{Statistical algorithms}

Specify rollout macro-aggregation, construction of paired init-cluster means,
deterministic cluster bootstrap, exhaustive sign-flip enumeration and Holm
family boundaries in sufficient detail to reproduce each interval and p-value.

\section{Supplementary results}
\label{app:supplementary}

The tables in this appendix are generated directly from canonical JSON/CSV
evidence by \texttt{build\_w1\_latex\_evidence.py}. Their source and output
hashes are recorded in \texttt{W1\_EVIDENCE\_TABLES\_COMPLETE.json}; displayed
precision does not replace the machine-readable values.

\subsection{All validation seeds and frozen-test checkpoints}

Table~\ref{tab:app-validation} exposes seed variability, capacity and selected
epoch for all 15 candidates. Ten runs select the last permitted epoch, so the
common budget may censor convergence even though it is equal across variants.

\begin{table}[p]
\centering\scriptsize
\caption{All validation runs used for frozen model selection. NLL, ADE and FDE are uncalibrated rollout-macro metrics: each is averaged within rollout and then equally across rollouts.}
\label{tab:app-validation}
\begin{tabular}{@{}llrrrrrr@{}}
\toprule
Variant & Seed & Epoch & Trainable & NLL & ADE & FDE & ms/sample \\
\midrule
B1 & 11 & 20 & 1,034,208 & 1.861 & 0.110 & 0.135 & 16.72 \\
B1 & 23 & 20 & 1,034,208 & 1.860 & 0.114 & 0.146 & 16.93 \\
B1 & 37 & 20 & 1,034,208 & 1.861 & 0.107 & 0.131 & 16.33 \\
B2-M & 11 & 18 & 77,600 & 1.997 & 0.914 & 1.626 & 18.32 \\
B2-M & 23 & 19 & 77,600 & 2.029 & 1.088 & 1.722 & 16.49 \\
B2-M & 37 & 20 & 77,600 & 2.026 & 1.053 & 1.807 & 18.05 \\
B2-D & 11 & 20 & 176,096 & 1.873 & 0.248 & 0.449 & 18.37 \\
B2-D & 23 & 19 & 176,096 & 1.873 & 0.251 & 0.444 & 17.64 \\
B2-D & 37 & 20 & 176,096 & 1.872 & 0.244 & 0.440 & 20.15 \\
T1 & 11 & 19 & 86,688 & 2.003 & 0.945 & 1.592 & 18.69 \\
T1 & 23 & 20 & 86,688 & 2.009 & 0.953 & 1.830 & 16.81 \\
T1 & 37 & 20 & 86,688 & 2.024 & 1.046 & 1.722 & 16.61 \\
T2 & 11 & 20 & 165,728 & 1.879 & 0.276 & 0.472 & 18.54 \\
T2 & 23 & 18 & 165,728 & 1.878 & 0.271 & 0.470 & 17.99 \\
T2 & 37 & 20 & 165,728 & 1.877 & 0.268 & 0.455 & 17.14 \\
\bottomrule
\end{tabular}
\end{table}


Table~\ref{tab:app-test} reports the one-shot test after model and representative
seed were frozen on validation. B0 was evaluated as a post-selection reporting
bridge and did not alter the selected B1 stack.

\begin{table}[h]
\centering\small
\caption{One-shot frozen-test results. Calibration was fitted on validation and was not used for architecture ranking.}
\label{tab:app-test}
\begin{tabular}{@{}llrrrrrr@{}}
\toprule
Variant & Seed & Val. rank & NLL & ADE & FDE & Cal. NLL & ms/sample \\
\midrule
B0 & -- & control & 2.171 & 1.299 & 2.685 & 0.582 & 17.62 \\
B1 & 37 & 1 & 1.857 & 0.106 & 0.129 & -2.069 & 18.51 \\
B2-D & 11 & 2 & 1.873 & 0.226 & 0.383 & -1.455 & 18.33 \\
T2 & 23 & 3 & 1.878 & 0.246 & 0.403 & -1.072 & 19.05 \\
T1 & 23 & 4 & 2.004 & 0.928 & 1.704 & -1.287 & 19.78 \\
B2-M & 37 & 5 & 2.024 & 1.047 & 1.730 & -1.351 & 18.74 \\
\bottomrule
\end{tabular}
\end{table}


\subsection{Complete corrected closed-loop contrasts}

Tables~\ref{tab:app-h3-contrasts} and~\ref{tab:app-h4-contrasts} contain every
primary metric contrast, including descriptive cluster intervals, exact
sign-flip p-values and Holm-adjusted values. No adjusted value is confirmatory;
this is consistent with five paired clusters and is not evidence of
equivalence.

\begin{table}[p]
\centering\scriptsize
\caption{All corrected R3 H3 metric contrasts. Effects are B1 minus B0; intervals are deterministic cluster-bootstrap descriptions and tests use five paired initialisation groups.}
\label{tab:app-h3-contrasts}
\begin{tabular}{@{}L{0.25\linewidth}llrL{0.15\linewidth}rrr@{}}
\toprule
Contrast & Metric & Style & Effect & 95\% CI & $p$ & $p_{\mathrm{Holm}}$ & $n$ \\
\midrule
B1--B0, fixed-aggressive & Completion (s) & assertive & -0.790 & [-1.410, -0.180] & 0.1250 & 1.0000 & 5 \\
B1--B0, fixed-aggressive & Separation (m) & assertive & +0.005 & [-0.011, +0.025] & 0.6875 & 1.0000 & 5 \\
B1--B0, fixed-aggressive & Completion (s) & reactive & -0.010 & [-0.090, +0.120] & 1.0000 & 1.0000 & 5 \\
B1--B0, fixed-aggressive & Separation (m) & reactive & -0.011 & [-0.031, +0.000] & 0.3750 & 1.0000 & 5 \\
B1--B0, fixed-medium & Completion (s) & assertive & +0.230 & [-0.630, +1.170] & 0.6250 & 1.0000 & 5 \\
B1--B0, fixed-medium & Separation (m) & assertive & +0.007 & [-0.000, +0.015] & 0.4375 & 1.0000 & 5 \\
B1--B0, fixed-medium & Completion (s) & reactive & -0.220 & [-0.430, -0.020] & 0.1875 & 1.0000 & 5 \\
B1--B0, fixed-medium & Separation (m) & reactive & -0.008 & [-0.023, +0.001] & 0.3750 & 1.0000 & 5 \\
B1--B0, fixed-conservative & Completion (s) & assertive & -0.520 & [-1.320, +0.090] & 0.4375 & 1.0000 & 5 \\
B1--B0, fixed-conservative & Separation (m) & assertive & -0.026 & [-0.075, +0.000] & 0.5000 & 1.0000 & 5 \\
B1--B0, fixed-conservative & Completion (s) & reactive & +0.490 & [+0.020, +0.960] & 0.2500 & 1.0000 & 5 \\
B1--B0, fixed-conservative & Separation (m) & reactive & -0.006 & [-0.016, +0.005] & 0.3750 & 1.0000 & 5 \\
B1--B0, adaptive & Completion (s) & assertive & +0.240 & [-0.090, +0.600] & 0.3750 & 1.0000 & 5 \\
B1--B0, adaptive & Separation (m) & assertive & -0.007 & [-0.018, +0.000] & 0.3125 & 1.0000 & 5 \\
B1--B0, adaptive & Completion (s) & reactive & -0.180 & [-0.660, +0.360] & 0.6250 & 1.0000 & 5 \\
B1--B0, adaptive & Separation (m) & reactive & +0.001 & [-0.016, +0.021] & 0.8125 & 1.0000 & 5 \\
\bottomrule
\end{tabular}
\end{table}


\begin{table}[p]
\centering\scriptsize
\caption{All corrected R3 H4 metric contrasts. Effects are adaptive minus fixed; Holm adjustment is within each registered three-comparator predictor--style family.}
\label{tab:app-h4-contrasts}
\begin{tabular}{@{}L{0.25\linewidth}llrL{0.15\linewidth}rrr@{}}
\toprule
Contrast & Metric & Style & Effect & 95\% CI & $p$ & $p_{\mathrm{Holm}}$ & $n$ \\
\midrule
B1: adaptive vs fixed-aggressive & Completion (s) & assertive & +0.530 & [+0.010, +1.200] & 0.2500 & 0.7500 & 5 \\
B1: adaptive vs fixed-aggressive & Separation (m) & assertive & -0.015 & [-0.037, +0.011] & 0.3750 & 1.0000 & 5 \\
B1: adaptive vs fixed-medium & Completion (s) & assertive & -0.440 & [-1.170, +0.040] & 0.3750 & 0.7500 & 5 \\
B1: adaptive vs fixed-medium & Separation (m) & assertive & -0.000 & [-0.014, +0.011] & 0.8125 & 1.0000 & 5 \\
B1: adaptive vs fixed-conservative & Completion (s) & assertive & +0.440 & [-0.050, +0.970] & 0.4375 & 0.7500 & 5 \\
B1: adaptive vs fixed-conservative & Separation (m) & assertive & -0.003 & [-0.017, +0.009] & 0.6875 & 1.0000 & 5 \\
B1: adaptive vs fixed-aggressive & Completion (s) & reactive & -0.140 & [-0.490, +0.100] & 0.6875 & 1.0000 & 5 \\
B1: adaptive vs fixed-aggressive & Separation (m) & reactive & +0.001 & [-0.000, +0.001] & 0.5000 & 1.0000 & 5 \\
B1: adaptive vs fixed-medium & Completion (s) & reactive & -0.140 & [-0.520, +0.130] & 0.6250 & 1.0000 & 5 \\
B1: adaptive vs fixed-medium & Separation (m) & reactive & +0.001 & [-0.001, +0.002] & 0.5000 & 1.0000 & 5 \\
B1: adaptive vs fixed-conservative & Completion (s) & reactive & -0.420 & [-0.860, +0.070] & 0.1875 & 0.5625 & 5 \\
B1: adaptive vs fixed-conservative & Separation (m) & reactive & -0.002 & [-0.013, +0.010] & 0.9375 & 1.0000 & 5 \\
B0: adaptive vs fixed-aggressive & Completion (s) & assertive & -0.500 & [-1.220, +0.030] & 0.3125 & 0.9375 & 5 \\
B0: adaptive vs fixed-aggressive & Separation (m) & assertive & -0.003 & [-0.023, +0.016] & 0.8750 & 1.0000 & 5 \\
B0: adaptive vs fixed-medium & Completion (s) & assertive & -0.450 & [-1.170, +0.120] & 0.5000 & 1.0000 & 5 \\
B0: adaptive vs fixed-medium & Separation (m) & assertive & +0.014 & [+0.001, +0.026] & 0.1875 & 0.5625 & 5 \\
B0: adaptive vs fixed-conservative & Completion (s) & assertive & -0.320 & [-0.980, +0.170] & 0.5625 & 1.0000 & 5 \\
B0: adaptive vs fixed-conservative & Separation (m) & assertive & -0.023 & [-0.077, +0.008] & 0.9375 & 1.0000 & 5 \\
B0: adaptive vs fixed-aggressive & Completion (s) & reactive & +0.030 & [-0.250, +0.410] & 0.9375 & 1.0000 & 5 \\
B0: adaptive vs fixed-aggressive & Separation (m) & reactive & -0.011 & [-0.051, +0.016] & 0.8125 & 1.0000 & 5 \\
B0: adaptive vs fixed-medium & Completion (s) & reactive & -0.180 & [-0.520, +0.160] & 0.5000 & 1.0000 & 5 \\
B0: adaptive vs fixed-medium & Separation (m) & reactive & -0.008 & [-0.030, +0.013] & 0.5625 & 1.0000 & 5 \\
B0: adaptive vs fixed-conservative & Completion (s) & reactive & +0.250 & [-0.210, +0.710] & 0.3750 & 1.0000 & 5 \\
B0: adaptive vs fixed-conservative & Separation (m) & reactive & -0.008 & [-0.031, +0.015] & 0.4375 & 1.0000 & 5 \\
\bottomrule
\end{tabular}
\end{table}


Table~\ref{tab:app-h4-dominance} applies the registered joint dominance rule to
the same contrasts. It is the source of the 3/12 H4 count and does not search
over comparators or geometry margins.

\begin{table}[h]
\centering\small
\caption{Registered H4 dominance decisions. Effects are adaptive minus fixed; lower completion and higher separation are preferred.}
\label{tab:app-h4-dominance}
\begin{tabular}{@{}llllrrl@{}}
\toprule
Predictor & Style & Fixed comparator & $\Delta t$ (s) & $\Delta d$ (m) & Binary guards & Decision \\
\midrule
B0 & assertive & fixed\_aggressive & -0.500 & -0.0032 & pass & Does not dominate \\
B0 & assertive & fixed\_conservative & -0.320 & -0.0228 & pass & Does not dominate \\
B0 & assertive & fixed\_medium & -0.450 & +0.0136 & pass & Dominates \\
B0 & reactive & fixed\_aggressive & +0.030 & -0.0107 & pass & Does not dominate \\
B0 & reactive & fixed\_conservative & +0.250 & -0.0081 & pass & Does not dominate \\
B0 & reactive & fixed\_medium & -0.180 & -0.0084 & pass & Does not dominate \\
B1 & assertive & fixed\_aggressive & +0.530 & -0.0151 & pass & Does not dominate \\
B1 & assertive & fixed\_conservative & +0.440 & -0.0033 & pass & Does not dominate \\
B1 & assertive & fixed\_medium & -0.440 & -0.0003 & pass & Does not dominate \\
B1 & reactive & fixed\_aggressive & -0.140 & +0.0005 & pass & Dominates \\
B1 & reactive & fixed\_conservative & -0.420 & -0.0019 & pass & Does not dominate \\
B1 & reactive & fixed\_medium & -0.140 & +0.0006 & pass & Dominates \\
\bottomrule
\end{tabular}
\end{table}


\subsection{Calibration and context ablations}

The following generated tables separate three distinct mechanism questions:
whether the Transformer uses its sequence, which B1 input drives aggregate
accuracy, and whether aggregate calibration represents the response-active
tail. An ablation delta establishes sensitivity to an input transformation,
not semantic understanding or closed-loop causality.

\begin{table}[h]
\centering\small
\caption{Post-selection sequence ablations on the frozen test set. Positive deltas mean ablation worsened the metric.}
\label{tab:app-sequence-ablation}
\begin{tabular}{@{}llrrr@{}}
\toprule
Variant & Ablation & $\Delta$NLL & $\Delta$ADE (m) & $\Delta$FDE (m) \\
\midrule
T1 & zero & +0.022 & -0.033 & +0.021 \\
T1 & shuffle & +0.085 & +0.144 & +0.284 \\
T2 & zero & +0.091 & +0.198 & +0.292 \\
T2 & shuffle & +0.149 & +0.350 & +0.428 \\
\bottomrule
\end{tabular}
\end{table}

\begin{table}[h]
\centering\small
\caption{B1 input-shuffle diagnostics across three deterministic mappings.}
\label{tab:app-b1-inputs}
\begin{tabular}{@{}llrrr@{}}
\toprule
Input & Mapping seed & $\Delta$NLL & $\Delta$ADE (m) & Active-tail $\Delta$ADE (m) \\
\midrule
Raster & 20260808 & +0.095 & +0.28883 & +1.907 \\
Target history & 20260808 & -0.000 & +0.00006 & +0.026 \\
Raster & 20260809 & +0.097 & +0.29293 & +1.923 \\
Target history & 20260809 & +0.000 & +0.00019 & +0.029 \\
Raster & 20260810 & +0.086 & +0.27053 & +1.890 \\
Target history & 20260810 & -0.000 & -0.00002 & +0.022 \\
\bottomrule
\end{tabular}
\end{table}

\begin{table}[h]
\centering\small
\caption{Aggregate and response-active calibration diagnostic for the frozen B0/B1 test stacks. The active tail has 15 windows from six rollouts and three initialisation groups.}
\label{tab:app-calibration-tail}
\begin{tabular}{@{}lllrrrr@{}}
\toprule
Subset & Stack & Windows & Rollouts & Init groups & Cal. NLL & Coverage MAE \\
\midrule
all & B0 & 315 & 20 & 5 & 0.582 & 0.407 \\
all & B1 & 315 & 20 & 5 & -2.069 & 0.077 \\
response\_active & B0 & 15 & 6 & 3 & 2.959 & 0.418 \\
response\_active & B1 & 15 & 6 & 3 & 8.573 & 0.443 \\
\bottomrule
\end{tabular}
\end{table}


\subsection{Collision-filtered sensitivity}

Six reactive training rollouts contained 253 raw callbacks. Canonical actor-pair
deduplication and contiguous-frame reconstruction produced 16 contact episodes
over 62 unique frames, primarily target--infrastructure contact. Exact sample-
clock alignment could not be recovered, so the sensitivity conservatively
removed all 162 usable windows from those six rollouts, regenerated train-only
normalisation and reran all 15 validation configurations. The five-model order
remained B1, B2-D, T2, T1, B2-M. This analysis never accessed test labels,
did not change the frozen primary dataset or selected checkpoint, and is
therefore a robustness check rather than a second selection experiment.

\subsection{Evidence not pooled with R3}

Day10--13 analyses use earlier control or geometry implementations and support
diagnostic mechanism statements only. They are absent from
Tables~\ref{tab:app-h3-contrasts}--\ref{tab:app-h4-dominance}, never enlarge
the five-cluster sample and never enter H3/H4 estimates.

\clearpage

\section{Claim-to-evidence audit}
\label{app:evidence}

The final evidence audit reports 14/14 hard checks passed, no unresolved result
identifiers and no source-hash mismatch. The canonical evidence package
contains 210 stable result IDs, eight tables and eight figures in both SVG and
PNG form. \TODO{Add the final Git commit and SHA-256 of the audit JSON used for
the submitted PDF.}

\subsection{Primary and secondary evidence boundary}

Primary evidence comprises the rollout-disjoint collection, three-seed
validation selection, frozen test, and the nominal/shifted formal closed-loop
matrices. Context ablations, calibration subsets, collision reconstruction and
Day13 filtered validation are diagnostic or sensitivity evidence. They support
mechanism and robustness claims but do not reopen model selection.

\subsection{Result provenance}

Every numerical sentence in the final paper should have an adjacent source
comment of the form:

\begin{verbatim}
% EVIDENCE: R_TEST_B1_MINUS_B0_ADE
% SOURCE: docs/paper/generated/paper_assets_v1/paper_key_results.csv
\end{verbatim}

The submitted archive should include the LaTeX source, bibliography, official
style assets, figures and a machine-readable claim-to-evidence table, while
excluding secrets, raw credentials and non-anonymised supplementary metadata.



\end{document}
